\scanpage{13}\mspage{\handwritten{12}}

$X$. De plus, on a la formule \inserted{\handwritten{du type $g^{\bullet}f_{\bullet}=f'_{\bullet}g'^{\bullet}$ ci-dessus, \unsure{impliquent|impliquant} la formule}} de projection habituelle. Ceci posé, les constructions bien connues en théorie des correspondances permettent de définir une catégorie $\underline{B}$ dont lds objets sont ceux de $C$, et les $\mathrm{Hom}$ étant donnés par
\[
\mathrm{Hom}_B(X,Y)=B(X\mathbin{x}Y)\ .
\]

On a un isomorphisme canonique $\underline{B}\handwritten{\simeq}\underline{B}^{o}$, qui est l'identité sur les objets et la symmétrie $B(X\mathbin{x}Y)\handwritten{\simeq}B(Y\mathbin{x}X)$ sur les flêches, et un foncteur évident \deleted{\unsure{$C\longrightarrow\underline B$}} $F_{\bullet}:C\handwritten{\to}\underline{B}$, qui est l'identité sur les objets et \unsure{donné} par les graphes sur les flêches. Composé avec la symmétrie précédente, on trouve un foncteur $F^{\bullet}:\deleted{\unsure{B(}}\,C\handwritten{\to}\underline{B}$. On vérifie enfin que le couple $(F_{\bullet},F^{\bullet})$ satisfait aux conditions envisagées dans l'alinéa 1 (ceci \unsure{aussi} devrait être formalisé au cas des correspondances générales ...). On vérifie enfin qu'il a la propriété universelle qu'il faut. Lorsqu'on se donne un couple $(F'_{\bullet},F'^{\bullet})$ à valeursdans une catégorie $B'$, il faut trouver une factorisation par $B$. Elle est toute trouvée sur les objets, il faut la définir sur les flêches, donc définir, pour $X,Y$ dans $C$, une application
\begingroup
\makeatletter\tagsleft@true\makeatother
\[
B(X\mathbin{x}Y)\handwritten{\longrightarrow}\mathrm{Hom}(X',Y')\ . \tag*{$(\ast)$}
\]
\endgroup
Mais si $Z$ est au\deleted{d}dessus de $X\mathbin{x}Y$, d'où $p:Z\handwritten{\to}X$ et $q:Z\handwritten{\to}Y$, on lui \unsure{associe} un élément du deuxième membre de $(\ast)$ comme le composé
\[
X'\xrightarrow{F'^{\bullet}(p)} Z'\xrightarrow{F'_{\bullet}(\deleted{\unsure{p}}q)}Y'\ .
\]

Il faut prouver que l'élément obtenu en remplaçant $Z$\inserted{\handwritten{\unsure{$Z_o$|$(=Z_0)$}}} par un $Z_1$ cobordant au dessus de $X$ \inserted{\handwritten{est le même}}. Or si $g:T\handwritten{\to}X\mathbin{x}Y$ réalise le cobordisme, on voit (regardant seulement $T$ et la décomposition $\mathrm{bord}(T)=Z_{\unsure{o|0}}\handwritten{\amalg}Z_1$) qu'il existe une variété $\handwritten{\overline{T}}$ sur $S^1$ (obtenue en ``doublant'' $T$) \deleted{et une application $T$ \unsure{sur}} trqnsversale au dessus d'un couple de points $s_o,\unsure{s_1}$ de $S^1$, tels que $Z_i$ soit l'image inverse \inserted{\handwritten{dans $\overline{T}$}} de $s_i$ ($i=0,1$); d'autre part,

% archive sheet number: 12
